\documentclass[10pt,a4paper]{amsart}
\usepackage[margin=19mm]{geometry}
\usepackage{amsmath,amssymb,mathtools}
\usepackage[scr=rsfs,cal=euler]{mathalfa}
\usepackage{xfrac}
\usepackage[unicode,hidelinks,bookmarks=false]{hyperref}
\newcommand{\ie}{i.e.\ }
\newcommand{\cf}{cf.\ }
\newcommand{\resp}{respectively\ }
\newcommand{\Subsection}{\subsection}
\providecommand{\nobreakdash}{\nobreak\hbox{-}}
\setlength{\emergencystretch}{2em}
\multlinegap0pt
\usepackage{amssymb}
\newtheorem{theorem}{Theorem}
\title{On the 3-rank of the class group of quadratic fields}
\author{Shi-Chao Chen, Chuan-Chuan Wu}
\date{Source: arXiv:2512.20023v1 (CC BY 4.0)}
\begin{document}
\maketitle

\noindent\emph{Source and licence.} On the 3-rank of the class group of quadratic fields. Version arXiv:2512.20023v1 (CC BY 4.0), distributed by arXiv under the Creative Commons Attribution 4.0 International licence, http://creativecommons.org/licenses/by/4.0/, as recorded in the arXiv licence metadata for this exact version and shown on the abstract page https://arxiv.org/abs/2512.20023v1. Full licence notice: \texttt{licenses/arxiv-2512.20023-CC-BY-4.0.txt}.\par\smallskip

\noindent\textit{Editorial scope.} Counted unit: the article's theorem be (source0.tex lines 158--165). The article's complete original proof of it is delivered with it (source0.tex lines 167--191, one \texttt{proof} environment of 25 lines). No mathematics is written, repaired or re-derived: the statement and the proof are the article's own lines, in the article's own notation, and no retained line is substituted or re-flowed.  Retained uncounted as the source-local context the proof needs: lines 138--149.  Nothing was omitted from any retained span, so the package carries no omission record.  No retained line contains a citation, so the wrapper carries no bibliography and no bibliography entry.  No retained line contains a figure environment, an image-inclusion command or drawing code, so no image is delivered and the package declares no asset dependency.  Editorial formatting only, in addition: an AMS \texttt{amsart} preamble that restates the article's own declarations and macros used by the retained lines, namely the environments \texttt{theorem} on separate counters, together with the standard packages those lines need.  The retained environments print in this document as Theorem 1 in order of appearance, whereas the article prints the same items with the numbering of its own full document.  The complete native source of the article as distributed in the arXiv e-print for this exact version is bundled unchanged as \texttt{sources/191527/source0.tex}.
\begin{theorem}[Nakagawa and Horie]\label{nh}
Suppose that the pair $(m,N)\in\mathbb{N}^2$ is good. Then we have


    \[
    \sum_{D \in S_{\lambda}(X, m, N)} 3^{r_{3}(D)} \sim c(\lambda)  \sum_{D \in S_{\lambda}(X, m, N)}1,
    \]
    where $c(\lambda)=\frac{4}{3}$ if $\lambda=1$, and $2$ if $\lambda=-1$. Moreover,
    \[\sum_{D \in S_{\lambda}(X, m, N)}1\sim \frac{3X}{\pi^2 \varphi(N)}\prod_{p|N}\frac{q}{1+p},\]
where $\varphi(\cdot)$ is the Euler's totient function, $q=4$ if $p=2$, and $q=p$ otherwise.

\end{theorem}

\begin{theorem}\label{be}
Suppose that the pair $(m,N)\in\mathbb{N}^2$ is good. Then
 for any $n \in \mathbb{N}$, we have
\begin{equation}\label{1}
    \liminf_{X \rightarrow \infty} \frac{\sharp\left\{D \in S_{\lambda}(X, m, N) \mid r_{3}(D) <n \right\}}{\sharp S_{\lambda}(X, m, N)} \ge\frac{3^n-c(\lambda) }{3^{n}-1},
\end{equation}
where $c(\lambda)=\frac{4}{3}$ if $\lambda=1$, and $2$ if $\lambda=-1$.
\end{theorem}

\begin{proof}
We observe that
\[
\begin{aligned}
 \sum_{D \in S_{\lambda}(X, m, N)} 3^{r_{3}(D)}
& =\sum_{\substack{D \in S_{\lambda}(X, m, N) \\
r_{3}(D)<n}} 3^{r_{3}(D)}+\sum_{\substack{D \in S_{\lambda}(X, m, N) \\
r_{3}(D) \geq n}} 3^{r_{3}(D)} \\
& \ge \sum_{\substack{D \in S_{\lambda}(X, m, N) \\
r_{3}(D)<n}} 1+3^n\sum_{\substack{D \in S_{\lambda}(X, m, N) \\
r_{3}(D) \geq n}} 1\\
&=\sum_{\substack{D \in S_{\lambda}(X, m, N) \\
r_{3}(D)<n}} 1+3^n\left(\sum_{\substack{D \in S_{\lambda}(X, m, N) }} 1-\sum_{\substack{D \in S_{\lambda}(X, m, N) \\
r_{3}(D) < n}} 1\right)\\
&=3^n\sum_{\substack{D \in S_{\lambda}(X, m, N) }} 1+(1-3^n)\sum_{\substack{D \in S_{\lambda}(X, m, N) \\
r_{3}(D) < n}} 1.
\end{aligned}
\]
Thus
\[ \sum_{\substack{D \in S_{\lambda}(X, m, N) \\
r_{3}(D) < n}} 1
\ge \frac{3^n}{3^n-1}\sum_{\substack{D \in S_{\lambda}(X, m, N) }} 1-\frac{1}{3^n-1}\sum_{D \in S_{\lambda}(X, m, N)} 3^{r_{3}(D)}.\]
Dividing $\sum_{\substack{D \in S_{\lambda}(X, m, N) }}1$ and applying Theorem \ref{nh}, we complete the proof.

\end{proof}

\end{document}
